\chapter{Start and Continue Daily Work}\label{ch:daily}
\section{Choose the result before the client}

A useful starting request names what should be different when the work is done. For
example, ``explain where these two cutflows first diverge'' defines an investigation;
``change the selection and show its effect'' additionally defines implementation and
scientific validation. Before acting, identify the owning repository or document, the
evidence available, and the result the reader will inspect. This determines whether the
work needs a short interactive exchange or a durable session \cite[rules
  1--3]{skill:coordination}.

For any task beyond a quick edit, start the agent in \code{scientific-mode}. It states
the goal back in its own words, picks the playbook that fits (a study or comparison,
figures and reports, or picking up, pausing and closing work) and copies the playbook's
steps into its todo list, so a skipped step is visible rather than forgotten \cite[Start;
  Playbooks]{skill:mode}.

Routine, mechanical, local work can remain lightweight. A spelling correction or a
bounded inspection does not need the full scaffolding of a remote study. Delegated,
unattended, cross-host, consequential, long-running, substantial, or recovery-worthy work
requires \code{BRIEF.md}, \code{STATE.md}, and \code{RESULT.md}. Add a machine-readable
\code{HANDOFF.json} when programmatic dependencies or cross-harness recovery require one
\cites[rule 3]{skill:coordination}[Handoff marker]{skill:jobs}.

The brief preserves the commitment: question, owner, write boundary, evidence,
verification, allowed actions, and stopping conditions. State records the actual
checkpoint, remaining work, and next safe action. The result states what was produced and
which claims remain limited. These records serve different points in the same task;
keeping state current makes a later handoff much easier \cite[Dispatch
  packet]{skill:delegation}.

For a substantial task, make the completion sentence concrete. ``The comparison report is
ready for my review when the inputs and cutflows have been checked, the report passes
internal reviews, and the exact PDF is available locally'' identifies an assessable
intermediate outcome. Human acceptance remains open. A code-only task can instead name
its implementation, checks, and reading guide; it acquires a report requirement only when
the task contract calls for one \cites[rule 4]{skill:coordination}[Verify and hand
  back]{skill:coding}.

\section{Identify the owner and the execution boundary}

Before choosing commands, distinguish four roles: the control surface, the agent runtime,
the owned checkout, and any execution target. Figure~\ref{fig:roles} shows their
relationship. A phone can steer an agent on a workstation whose local worktree contains
the source while an authorized job runs on another host. The phone's location does not
relocate that source or change the job's owner \cite[GUI, terminal and cross-client
  sessions]{skill:workflow-gui}.

\begin{figure}[htbp]
  \centering
  \begin{tikzpicture}[node distance=10mm]
    \node[work,text width=110mm] (ui) {\textbf{Control surface}\\phone or laptop;
      choose the intended session};
    \node[work,text width=110mm,below=of ui] (runtime) {\textbf{Agent runtime and owned worktree}\\
      observed host, repository, branch, writing owner};
    \node[check,text width=110mm,below=of runtime] (run) {\textbf{Execution target, when needed}\\
      capable environment and explicitly authorized job};
    \draw[flow] (ui)--node[right,font=\footnotesize]{steer}(runtime);
    \draw[flow] (runtime)--node[right,font=\footnotesize]{run / inspect}(run);
  \end{tikzpicture}
  \caption{Control, source ownership, and execution are separate roles. This is a
    routing model, not evidence of a particular live connection.}
  \label{fig:roles}
\end{figure}

Inspect the literal host, Git state, worktree, and active session before editing. An
unfamiliar dirty path needs ownership evidence. Preserve it while resolving the
uncertainty; do not reset it to make the checkout look convenient. On the same machine,
use an isolated worktree when report or software edits would overlap another task
\cite[Start a writing task; Commit and hand off]{skill:git}.

\hook{daily-continuation}

\section{Keep the work loop small and the ending explicit}

Within the owned boundary, alternate a coherent change with the evidence needed to assess
it. Remote execution belongs inside that loop whenever the check needs remote data or
software. Figure~\ref{fig:loop} shows the dependency without forcing every task through a
report or putting all validation before execution.

\begin{figure}[htbp]
  \centering
  \begin{tikzpicture}[node distance=14mm]
    \node[work,text width=45mm] (do) {\textbf{Implement or investigate}\\
      bounded question and owned change};
    \node[check,text width=45mm,right=18mm of do] (test) {\textbf{Run and inspect}\\
      local or authorized remote evidence};
    \node[result,text width=108mm,below=22mm of do.south east,anchor=north] (hand)
    {\textbf{Hand back the promised result}\\applicable reviews, artifact,
      remaining decisions, and publication state};
    \draw[flow] (do)--(test);
    \draw[flow] (test.north) to[bend right=35] node[above,font=\footnotesize]
    {failed check: revise} (do.north);
    \draw[flow] (test.south)--(hand.north east);
  \end{tikzpicture}
  \caption{A shared work loop with conditional outputs. Source publication can
    occur at authorized durability checkpoints; report review can return to revision
    \cites[Verify and hand back]{skill:coding}[Review and release]{skill:report-production}
    [Commit and hand off]{skill:git}.}
  \label{fig:loop}
\end{figure}

Close only the gates named by the task. Code checks, source publication, PDF
availability, external delivery, notification, and human acceptance are separately
observable states. Record an absent requirement as not applicable and an unmet
requirement as pending. A substantial report remains \texttt{awaiting human review} after
internal checks and requested delivery until its owner accepts, supersedes, or closes it.
A version that an agent replaces before the owner has reviewed it is marked
\texttt{replaced before review}; \texttt{superseded} is reserved for the owner's own
review or decision \cite[Review and release]{skill:report-production}.
